%!TEX root = main.tex
%!TEX spellcheck = en-US

In~\cite{DBLP:journals/logcom/MarinPPS25} a pure sequent calculus system ($\LE$) equivalent to Prawitz's ecumenical natrural deduction system $\NE$ was proposed\footnote{See \cite{DBLP:conf/dali/MarinPPS20,DBLP:journals/synthese/PimentelPP21,DBLP:conf/wollic/MarinPPS21,DBLP:conf/calco/PimentelP23,Pereira2024} for further discussion on ecumenism, harmony and purity of ecumenical systems.}. The motivation departs from the observation that ``Although being a powerful tool for describing proof-theoretical properties of Prawitz's $\NE$ is not satisfactory as a logical system since it is not pure''. 

The strategy in building a pure system is to (i) work with ``a new notion of polarities for ecumenical formulas''; 
and (ii) use Girard's notion of stoup \cite{DBLP:journals/mscs/Girard91} to control rule applications by separating contexts into bins. In their system, sequents have the form $\Gamma \Rightarrow \Delta;\Sigma$ where $\Gamma, \Delta, \Sigma$ represent multisets of formulas, and the \textit{stoup} $\Sigma$, is limited to containing exactly one formula. Two structural rules, \textit{Store} and \textit{Dereliction}, control how formulas move in and out of the \textit{stoup} area:

    \begin{center}
    \begin{minipage}[b]{.3 \textwidth}
        \begin{prooftree}
          \AxiomC{$\Gamma \Rightarrow \Delta,P; P$}
          \UnaryInfC{$\Gamma \Rightarrow \Delta,P;A$}   
        \end{prooftree}
    \end{minipage}
    \begin{minipage}[b]{.3 \textwidth}
        \begin{prooftree}
        \AxiomC{$\Gamma \Rightarrow \Delta, N;\bot$}
        \UnaryInfC{$\Gamma \Rightarrow \Delta; N$}
        \end{prooftree}
    \end{minipage}
    \end{center}

Our proposal is to develop a game semantics with respect to the pure ecumenical system described above with two players, $\I$  and $\Y$.

{\bf TODO:}
\begin{itemize}
\item Define positive and negative formulas. In particular, should we have negation? Add the discussion about 2 negations and how it does not affect {\em provability}, but it does matter for the definition of semantics and game steps.
\item Question: should we have int. and classical atoms? In my opinion we should not! Hence, if you want to win the game $\neg\neg p\to p$ with an intuitionistic negation, then the implication {\em must be}  classical and the cost is 100. Otherwise one could have an intuitionistic implication and the cost would be 0001...
\item I've commented weakening, since it can be simulated using dereliction and negation. Is this okay?
\end{itemize}


\begin{figure}[htp]
{\sc Intuitionistic and neutral Rules}
$$
\infer[{\wedge L}]{\Seq{\Gamma, A \wedge B}{ \Pi} {\Delta}}{\Seq{\Gamma,
A, B}{ \Pi}{\Delta}} 
\quad
\infer[{\wedge R}]{\Seq{\Gamma}
{A \wedge B}{\Delta}}{\deduce{\Seq{\Gamma} {A}{\Delta}}{} &\deduce{\Seq{\Gamma}{B}{\Delta}}{}} 
$$
$$
\infer[{\vee_i L}]{\Seq{\Gamma, A \vee_i B}{ \Pi}{\Delta}
}{\deduce{\Seq{\Gamma, A}{ \Pi}{\Delta}}{} &\deduce{\Seq{\Gamma,B}{\Pi}{\Delta}}{}} 
\quad
\infer[{\vee_i R_j}]{\Seq{\Gamma}{ A_1\vee_i A_2}{\Delta}}{\Seq{\Gamma}{ A_j}{\Delta}} 
$$
$$
\infer[{\iimp L}]{\Seq{\Gamma, A\iimp B}{
\Pi}{\Delta}}
{\Seq{\Gamma, A\iimp B}{ A}{\Delta} & \Seq{\Gamma , B} 
{\Pi}\Delta} 
\quad
\infer[{\iimp R}]{\Seq{\Gamma}
{A\iimp B}{\Delta}}{\Seq{\Gamma,A}{B}{\Delta}}
$$
$$
\infer[{\neg L}]{\Seq{\Gamma,\neg A}{C}{\Delta}
}{\Seq{\Gamma, \neg A}{  A}{\Delta}}
\quad
\infer[{\neg R}]{\Seq{\Gamma} 
{\neg A}{\Delta}}{\Seq{\Gamma,A}{\bot}{\Delta}}
$$
{\sc Classical Rules}
$$
\infer[\cimp L]{\Seq{\Gamma, A \rightarrow_c B}{C}{\Delta} }{\Seq{\Gamma, A \rightarrow_c B}{A}{\Delta}&\Seq{\Gamma,  B}{\bot}{\Delta}}
\quad
\infer[\cimp R]{\Seq{\Gamma} {C}{ A\rightarrow_c B,\Delta}}{\Seq{\Gamma, A}{  \bot}{B,\Delta}} 
$$
$$
\infer[\vee_cL]{\Seq{\Gamma,  A \vee_c B}{C}{\Delta}}{\Seq{\Gamma,A}{ \bot}{\Delta} & \Seq{\Gamma,B}{ \bot}{\Delta}} 
\quad
\infer[\vee_cR]{\Seq{\Gamma}{C}{A \vee_c B, \Delta}}{\Seq{\Gamma}{ \bot}{A,B,\Delta}}
$$
$$
\infer[L_c]{\Seq{\Gamma, p_c}{C}{\Delta}}{\Seq{\Gamma,p_i}{ \bot}{\Delta}}
\quad
\infer[R_c]{\Seq{\Gamma}{C}{p_c,\Delta}}{\Seq{\Gamma}{ \bot}{p_i,\Delta}} 
$$
{\sc Initial, Decision and Structural Rules}
$$
\infer[\init_i]{\Seq{\Gamma,A}{A}{\Delta}}{} 
\qquad
\infer[\init_c]{\Seq{\Gamma,A}{\Pi}{A,\Delta}}{} 
$$
$$
\infer[\D]{\Seq{\Gamma}{C}{P,\Delta}}{\Seq{\Gamma}{ P}{P,\Delta}} 
\qquad
\infer[\store]{\Seq{\Gamma}{N}{\Delta}}{\Seq{\Gamma}{ \bot}{N,\Delta}}
%\qquad
%\infer[\W]{\Seq{\Gamma}{A}{\Delta}}{\Seq{\Gamma}{ \bot}{\Delta}} 
$$
{\sc Cut Rules}
$$
\begin{array}{c@{\qquad}c}
\infer[\cut_P]{\Seq{\Gamma}{ \Pi}{\Delta}}{\Seq{\Gamma}{P}{\Delta} & \Seq{P,\Gamma}{ \Pi}{\Delta}} 
&
\infer[\cut_N]{\Seq{\Gamma}{ \Pi}{\Delta}}{\Seq{\Gamma}{ \Pi^*}{\Delta, N} & \Seq{N,\Gamma}{ \Pi}{\Delta}} 
\end{array}
$$
\caption{Ecumenical pure  system $\LE$. $N$ is negative and $P$ is positive; $p$ is atomic.  $\Pi^*$ is either $\bot$ or a $P\in\Delta$.}\label{fig:labEK}
\end{figure}


\subsection{Ecumenical birelational models}\label{sec:bi}
In~\cite{luiz17}, an ecumenical Kripke semantics was introduced in which the classical connectives are interpreted via a double-negation translation. Their semantic clauses thus inherit the dependence on accessible worlds characteristic of intuitionistic negation. Here, we propose an alternative ecumenical semantics in which classical behaviour is captured through a classical notion of negation. 

The resulting clauses for the classical connectives are \emph{local}: they refer to the evaluation of their immediate subformulas at the current world, without introducing quantification over accessible worlds. %This way a classical version of the game would coincide with provability in classical logics!

\elaine{I am not sure how to define this!}

\begin{definition}\label{def:ekripke} A {\em birelational Kripke model} is a quadruple  
$\m=(W,\leq,R,V)$ where $(W,R,V)$ is a Kripke model
such that $\wld$ is partially ordered with order $\leq$, $R\subset W\times W$ is a binary relation, the 
satisfaction function $V:\langle W,\leq\rangle\rightarrow  \langle2^\Prop,\subseteq\rangle$ is monotone and:

\noindent
F1. For all worlds $w,v,v'$, if $wRv$ and $v \leq v'$, there is a $w'$ such that $w \leq w'$ and $w'Rv'$;

\noindent
F2. For all worlds $w', w, v$, if $w \leq w'$ and $wRv$, there is a $v'$ such that $w'Rv'$ and $v \leq v'$.

An {\em ecumenical modal Kripke model} is a birelational Kripke model such that truth of an ecumenical formula at a point $w$ is the
smallest relation $\modelse$ satisfying 

\noindent
$
\begin{array}{l@{\quad}c@{\quad}l}
\m,w\modelse p & \mbox{ iff } & p\in \val(w);\\
\m,w\modelse A\wedge B & \mbox{ iff } & \m,w\modelse A \mbox{ and } \m,w\modelse B;\\
\m,w\modelse A\vee_i B & \mbox{ iff } & \m,w\modelse A \mbox{ or } \m,w\modelse B;\\
\m,w\modelse A\iimp B & \mbox{ iff } & \text{for all } v \text{ such that  }w \leq v,   
\m,v\modelse A \mbox{ implies }  \m,v\modelse B;\\
%\m,w\modelse \neg A & \mbox{ iff } & \text{for all } v \text{ such that  }w \leq v, \m,v\not\modelse A;\\
\m,w \modelse \bot & & \mbox{never holds};\\
\m,w\modelse A & \mbox{ iff } & \m,w\modelsec A \mbox{ if $A$ is classical};\\
\m,w\modelsec A & \mbox{ iff } & \m,w\modelse A \mbox{ if $A$ is neutral/intuitionistic};\\
%\m,w\modelse p_c  & \mbox{ iff } & \m,w\modelse \neg(\neg p_i);\\
\m,w\modelsec \neg A & \mbox{ iff } &   \magenta{\m,w\not\modelsec A};\\
\m,w\modelsec A\vee_c B & \mbox{ iff } & \magenta{\m,w\modelsec A \mbox{ or } \m,w\modelsec B};\\
\m,w\modelsec A\cimp B & \mbox{ iff } & \magenta{\m,w\not\modelsec  A \mbox{ or } \m,w\modelsec B}.
\end{array}
$

We say that a formula $A$ is {\em valid} in a model $\m=(W,\leq,R,V)$ if for all $w\in  W$ we have $w\modelse A$. A formula $A$ is {\em valid} in a frame $\langle W,\leq,R\rangle$ if, for all valuations $V$, $A$ is valid in the model $(W,\leq,R,V)$. Finally, we say a formula is {\em valid}, if it is valid in all frames.
\end{definition}
Since, restricted to intuitionistic and neutral connectives, $\modelse$ is the usual birelational interpretation $\models$ for $\IK$~\cite{plotkin:stirling:86,Sim94}, and since the classical connectives are interpreted 
via the neutral ones using the double-negation translation, an ecumenical modal Kripke model coincides with the standard birelational Kripke model for intuitionistic modal logic $\IK$. Hence the following result easily holds from the similar result for $\IK$.
\begin{theorem}[\cite{DBLP:conf/dali/MarinPPS20}]\label{thm:scLEci} The system
$\labEK$ is sound and complete w.r.t.~the ecumenical modal Kripke semantics, that is, $\vdash_{\labEK} x: A$ iff $\modelse A$.
\end{theorem}
We end this section with a small note on the relationship between the semantics and the dynamics of proofs. 
On a bottom-up reading of proofs, the $\store$ rule is a delay on applying rules over classical connectives. It corresponds to  moving the formula up w.r.t. $\leq$ in the birelational semantics. The rule $\Box R$, on the other hand, slides the formula to a fresh new world, related to the former one through the relation $R$. Finally, rule $\neg R$ moves up the formula w.r.t. $\leq$. 
